\chapter{Database Tables}
The dictionary documents JPA attributes and Java types, with explicit column names where declared. PK means primary key, FK relationship, NN non-null, and UQ unique. All IDs use identity-generated Long values; enums are stored as strings. Initialization values describe application defaults, not SQL DEFAULT clauses. Physical naming follows ORM configuration.

\section{Users}
\begingroup\footnotesize\setstretch{1.0}\setlength{\parskip}{0pt}\renewcommand{\arraystretch}{1.15}
\begin{longtable}{>{\raggedright\arraybackslash}p{4.5cm}>{\raggedright\arraybackslash}p{2.2cm}>{\raggedright\arraybackslash}p{1.4cm}>{\raggedright\arraybackslash}p{5.7cm}}
\caption{Data dictionary: \code{users}.}\\\toprule Attribute & Type & Mapping & Purpose / initialization\\\midrule\endfirsthead
\toprule Attribute & Type & Mapping & Purpose / initialization\\\midrule\endhead
\code{id} & Long & PK & Identity-generated account identifier.\\
\code{username} & String & NN, UQ & Up to 100 characters.\\
\code{email} & String & NN, UQ & Up to 150 characters; normalized during registration.\\
\code{password} & String & NN & Encoded password; authentication uses BCrypt matching.\\
\code{active} & Boolean & NN & Initialized to true; false blocks account access.\\
\code{role} & UserRole & NN & String enum, length 20; initialized to USER when absent.\\
\code{createdAt} & LocalDate\allowbreak{}Time & NN & Set on creation when absent.\\\bottomrule
\end{longtable}\endgroup

\noindent\textbf{Relationships:} one user can have many login records, anomalies, and analysis-history records. The audit entity stores a username string and does not declare a relationship to \code{users}. UserRole values are ADMIN and USER.

\section{Login Events}
\begingroup\footnotesize\setstretch{1.0}\setlength{\parskip}{0pt}\renewcommand{\arraystretch}{1.15}
\begin{longtable}{>{\raggedright\arraybackslash}p{4.5cm}>{\raggedright\arraybackslash}p{2.2cm}>{\raggedright\arraybackslash}p{1.4cm}>{\raggedright\arraybackslash}p{5.7cm}}
\caption{Data dictionary: \code{user_login_logs}.}\\\toprule Attribute / column & Type & Mapping & Purpose\\\midrule\endfirsthead
\toprule Attribute / column & Type & Mapping & Purpose\\\midrule\endhead
\code{id} & Long & PK & Login-event identifier.\\
\code{user} / \code{user_id} & User & FK, NN & Required many-to-one, lazy user reference.\\
\code{loginTime} / \code{login_time} & LocalDate\allowbreak{}Time & NN & Event timestamp; initialized when absent.\\
\code{ipAddress} / \code{ip_address} & String & NN & Up to 100 characters.\\
\code{location} & String & Optional & Up to 100 characters; event-supplied context.\\
\code{device} & String & Optional & Up to 200 characters; device/user-agent string.\\
\code{status} & LoginStatus & NN & String enum, length 20: SUCCESS or FAILED.\\\bottomrule
\end{longtable}\endgroup
Indexes are declared on \code{user_id}, \code{login_time}, \code{status}, and \code{ip_address}. Their entity names are \code{idx_login_username}, \code{idx_login_time}, \code{idx_login_status}, and \code{idx_login_ip}, respectively. Despite its name, \code{idx_login_username} indexes the user identifier.

\section{Anomalies}
\begingroup\footnotesize\setstretch{1.0}\setlength{\parskip}{0pt}\renewcommand{\arraystretch}{1.15}
\begin{longtable}{>{\raggedright\arraybackslash}p{4.5cm}>{\raggedright\arraybackslash}p{2.2cm}>{\raggedright\arraybackslash}p{1.4cm}>{\raggedright\arraybackslash}p{5.7cm}}
\caption{Data dictionary: \code{anomalies}.}\\\toprule Attribute / column & Type & Mapping & Purpose\\\midrule\endfirsthead
\toprule Attribute / column & Type & Mapping & Purpose\\\midrule\endhead
\code{id} & Long & PK & Anomaly identifier.\\
\code{user} / \code{user_id} & User & FK, NN & Required many-to-one account reference.\\
\code{loginLog} / \code{login_log_id} & LoginLog & FK, NN & Required triggering-event reference.\\
\code{anomalyType} & AnomalyType & NN & One of the six rule types.\\
\code{description} & String & NN & Human-readable trigger explanation.\\
\code{riskScore} & Integer & NN & Fixed weight assigned by the matching rule.\\
\code{severity} & Severity & NN & LOW, MEDIUM, HIGH, or CRITICAL.\\
\code{detectedAt} & LocalDate\allowbreak{}Time & NN & Detection timestamp initialized when saved.\\\bottomrule
\end{longtable}\endgroup

The six types are \code{FAILED_LOGIN}, \code{UNUSUAL_LOGIN_TIME}, \code{INACTIVE_USER_LOGIN}, \code{NEW_IP}, \code{NEW_DEVICE}, and \code{BRUTE_FORCE}. The service checks for an existing user/type/login combination before insertion; the entity does not declare a corresponding composite uniqueness constraint.

\section{AI Analysis History}
\begingroup\footnotesize\setstretch{1.0}\setlength{\parskip}{0pt}\renewcommand{\arraystretch}{1.15}
\begin{longtable}{>{\raggedright\arraybackslash}p{4.5cm}>{\raggedright\arraybackslash}p{2.2cm}>{\raggedright\arraybackslash}p{1.4cm}>{\raggedright\arraybackslash}p{5.7cm}}
\caption{Data dictionary: \code{ai_analysis_history}.}\\\toprule Attribute / column & Type & Mapping & Purpose\\\midrule\endfirsthead
\toprule Attribute / column & Type & Mapping & Purpose\\\midrule\endhead
\code{id} & Long & PK & Analysis record identifier.\\
\code{anomaly} / \code{anomaly_id} & Anomaly & FK, optional & Lazy link to the analyzed anomaly.\\
\code{user} / \code{user_id} & User & FK, optional & Lazy link to the associated user.\\
\code{cachedAnomalyId} & Long & Optional & Column \code{cached_anomaly_id}; retained identifier, not FK.\\
\code{cachedUserId} & Long & Optional & Column \code{cached_user_id}; retained identifier, not FK.\\
\code{username} & String & Optional & Up to 100 characters.\\
\code{analysisType} & String & NN & Up to 50 characters.\\
\code{summary} & String / TEXT & Optional & Generated or fallback overview.\\
\code{riskExplanation} & String / TEXT & Optional & Contextual risk explanation.\\
\code{indicatorsJson} & String / TEXT & Optional & Serialized indicators.\\
\code{recommendationsJson} & String / TEXT & Optional & Serialized recommendations.\\
\code{confidence} & String & Optional & Up to 20 characters; advisory label.\\
\code{threatLevel} & String & Optional & Up to 20 characters.\\
\code{model} & String & Optional & Up to 100 characters.\\
\code{isAdvisory} & Boolean & NN & Initialized to true.\\
\code{createdAt} & LocalDate\allowbreak{}Time & NN & Analysis creation timestamp.\\\bottomrule
\end{longtable}\endgroup

The JSON-labelled fields are mapped as TEXT, not as PostgreSQL JSONB in these entity definitions. Optional links and cached identifiers retain contextual information; the history DTO and service determine what is returned to the caller.

\section{Audit Records}
\begingroup\footnotesize\setstretch{1.0}\setlength{\parskip}{0pt}\renewcommand{\arraystretch}{1.15}
\begin{longtable}{>{\raggedright\arraybackslash}p{4.5cm}>{\raggedright\arraybackslash}p{2.2cm}>{\raggedright\arraybackslash}p{1.4cm}>{\raggedright\arraybackslash}p{5.7cm}}
\caption{Data dictionary: \code{audit_logs}.}\\\toprule Attribute & Type & Mapping & Purpose\\\midrule\endfirsthead
\toprule Attribute & Type & Mapping & Purpose\\\midrule\endhead
\code{id} & Long & PK & Audit record identifier.\\
\code{username} & String & Optional & Actor/account text; no declared user FK.\\
\code{action} & String & NN & Security or application action name.\\
\code{details} & String & Optional & Description, up to 1000 characters.\\
\code{ipAddress} & String & Optional & Recorded request/action source.\\
\code{timestamp} & LocalDate\allowbreak{}Time & NN & Initialized when absent.\\\bottomrule
\end{longtable}\endgroup
The audit table persists selected events. The entity does not establish append-only enforcement, hash chaining, or a database-trigger-based complete change log.
